% 太阳系（综述）
% license CCBYSA3
% type Wiki

本文根据 CC-BY-SA 协议转载翻译自维基百科\href{https://en.wikipedia.org/wiki/Solar_System}{相关文章}。

\begin{figure}[ht]
\centering
\includegraphics[width=8cm]{./figures/5a727e2db62d7a0e.png}
\caption{太阳、行星、卫星与矮行星\(^\text{[a]}\)（真实色彩，比例为尺寸等比，但距离不按比例）} \label{fig_Solar_1}
\end{figure}
太阳系（The Solar System）由太阳及其所环绕的天体组成。该名称来源于“Sōl”，即太阳的拉丁名称。大约在 46 亿年前，一个分子云中密度较高的区域发生坍缩，形成了太阳以及一个原行星盘（protoplanetary disc），环绕天体便由此汇聚而成。太阳核心内部的氢融合为氦，释放出能量，这些能量主要通过其外层光球（photosphere）向外辐射，从而在整个系统中形成一个由内而外逐渐降低的温度梯度。太阳系中超过 99.86\% 的质量集中在太阳内部。

在绕太阳运行的天体中，质量最大的为八大行星。按与太阳距离递增的顺序，最靠近太阳的四颗类地行星（terrestrial planets）分别是—— 水星（Mercury）、金星（Venus）、地球（Earth）和火星（Mars）。它们构成了太阳系的内行星系统（inner Solar System）。其中，地球与火星是太阳系中仅有的两颗位于太阳宜居带（habitable zone）内的行星——在此范围内，行星表面可存在液态水。

在距离太阳约 5 个天文单位（AU）之外、越过“霜线”（frost line）的位置，分布着四颗更为庞大的行星：两颗气态巨行星（gas giants）—— 木星（Jupiter）与土星（Saturn），以及两颗冰巨行星（ice giants）—— 天王星（Uranus）与 海王星（Neptune）。它们共同构成了太阳系的外行星系统（outer Solar System）。其中，木星与土星占据了太阳系除恒星外几乎 90\% 的全部质量。

太阳系中还存在着数量庞大但质量较小的天体。天文学家普遍认为，太阳系至少拥有九颗矮行星（dwarf planets），分别是：谷神星（Ceres）、奥库斯（Orcus）、冥王星（Pluto）、哈梅亚（Haumea）、夸奥尔（Quaoar）、鸟神星（Makemake）、共工（Gonggong）、阋神星（Eris）与赛德娜（Sedna）。在这些天体中，有六颗行星、七颗矮行星及若干其他天体拥有自然卫星（natural satellites），即通常所称的“月亮（moons）”。这些卫星的大小范围极广：从如地球之月那样接近矮行星尺度的巨型卫星，到仅有极小质量的“月粒（moonlets）”。此外，太阳系中还存在大量小型天体（small Solar System bodies），包括小行星（asteroids）、彗星（comets）、半人马天体（centaurs）、流星体（meteoroids）以及星际尘埃云（interplanetary dust clouds）。其中一部分天体分布在小行星带（asteroid belt）（位于火星与木星轨道之间）以及柯伊伯带（Kuiper belt）（位于海王星轨道之外）。

在这些天体之间，存在着充满尘埃与粒子的行星际介质（interplanetary medium）。太阳不断释放的带电粒子形成太阳风（solar wind），它向外扩散，构成一个被称为日球层（heliosphere）的巨大空间。在距离太阳约 70–90 个天文单位（AU）处，太阳风受到星际介质（interstellar medium）的阻挡，形成边界——日球顶（heliopause）。这一区域标志着太阳系与星际空间的分界。更远处、约 2,000 个天文单位（AU）以外，是太阳系最外层的假想区域——奥尔特云（Oort cloud）。它被认为是长周期彗星（long-period comets）的来源，其范围可能延伸至太阳系的引力界限（Hill sphere）边缘，约178,000–227,000 AU（2.81–3.59 光年），在此处太阳系的引力势与银河系的引力势达到平衡。目前，太阳系正穿越一个被称为本地星际云（Local Cloud）的星际介质区域。距离太阳系最近的恒星是比邻星（Proxima Centauri），距离约269,000 AU（4.25 光年）。太阳系与比邻星都位于本地泡（Local Bubble）内，这是银河系中一个直径约 1,000 光年（ly）的相对稀薄的区域。
\subsection{定义}
太阳系包括太阳以及所有受其引力约束并围绕其运行的天体。\(^\text{[15][16][17]}\)

国际天文学联合会（IAU）将太阳系描述为：由太阳引力所束缚的一切天体，包括太阳本身、八大行星，以及环绕太阳运行的其他天体。\(^\text{[11]}\)美国国家航空航天局（NASA）则将太阳系定义为一个行星系统（planetary system），其中包括太阳及所有环绕它运行的天体。\(^\text{[12]}\)

当“solar system”一词不作为专有名词、且首字母不大写时，它可以指太阳系本身，也可泛指任何类似太阳系的行星系统。\(^\text{[15]}\)
\subsection{形成与演化}
\subsubsection{过去}
\begin{figure}[ht]
\centering
\includegraphics[width=8cm]{./figures/0dc12bd0e5f767c3.png}
\caption{} \label{fig_Solar_2}
\end{figure}
太阳系形成于至少约46.68 亿年前，源自一片巨大分子云中某个区域的引力坍缩。\(^\text{[b]}\)这片最初的云可能横跨数光年，并极有可能同时孕育了多颗恒星。\(^\text{[19]}\)如同其他分子云一样，它的主要成分是氢气，其次是氦气，并含有少量由前几代恒星通过核聚变产生的较重元素。\(^\text{[20]}\)

当原太阳星云\(^\text{[20]}\)开始坍缩时，角动量守恒使其旋转速度不断加快。中心区域由于聚集了大部分质量，变得比周围更炽热、更致密。\(^\text{[16]}\)随着收缩的星云自转加速，它逐渐扁平化为一个直径约 200 天文单位（AU）的原行星盘\(^\text{[19][21]}\)，其中心形成了一个炽热而致密的原恒星。\(^\text{[22][23]}\)行星由盘中的物质通过吸积（accretion）过程形成，\(^\text{[24]}\)其中尘埃和气体相互引力吸引，逐渐聚合成更大的天体。在早期太阳系中，可能存在数百个原行星（protoplanets），但它们大多互相合并、被摧毁或被抛射出系统，最终仅留下行星、矮行星以及少量残余小天体。\(^\text{[25][26]}\)

在靠近太阳的内太阳系，即霜线（frost line）以内、甚至烟尘线（soot line）以内，\(^\text{[27]}\)由于高温使除金属与硅酸盐外的物质难以以固体形式存在，因此这一带形成的行星主要由岩石构成，即：水星、金星、地球与火星。由于这些耐高温的物质仅占太阳星云极小的比例，类地行星因此无法增长得太大。\(^\text{[25]}\)

四颗巨行星——木星、土星、天王星与海王星——在霜线之外形成，该区域位于火星与木星轨道之间，温度足够低，使挥发性冰类化合物可以保持固态。这些冰的丰度远高于类地行星形成所用的金属与硅酸盐，这使得巨行星能长得足够大，从而捕获大量的氢与氦 —— 这两种元素是宇宙中最轻、最丰富的。\(^\text{[25]}\)未能成为行星的残余碎片则聚集在小行星带、柯伊伯带与奥尔特云等区域。\(^\text{[25]}\)

约五千万年后，原恒星中心的氢气压与密度足够高，开始发生热核聚变反应。\(^\text{[28]}\)随着核心中氦的积累，太阳逐渐变得更亮；\(^\text{[29]}\)在主序星生命早期，其亮度仅为如今的约 70\%。\(^\text{[30]}\)温度、反应速率、压力与密度不断上升，直至达到静水平衡（hydrostatic equilibrium）：即热压与引力达到平衡。此时，太阳成为一颗主序星。\(^\text{[31]}\)太阳风从太阳表面喷射，形成了日球层（heliosphere），并将原行星盘中剩余的气体与尘埃吹散至星际空间。\(^\text{[29]}\)

在原行星盘消散后，尼斯模型（Nice model）提出：气体巨行星与微行星之间的引力相互作用导致行星迁移至不同轨道。这引发了整个系统的动力学不稳定，使微行星被散射，最终将气体巨行星带入如今的位置。在这一时期，大调车假说（Grand Tack Hypothesis）认为：木星的一次最终向内迁移扰乱了小行星带的大部分物质，从而引发了晚期重轰炸（Late Heavy Bombardment）。\(^\text{[32][33]}\)
\subsubsection{现在与未来}
太阳系目前处于一个相对稳定、缓慢演化的状态，所有天体都沿着受引力束缚的轨道绕太阳运行。\(^\text{[34]}\)尽管太阳系在数十亿年来一直较为稳定，但从技术角度上看，它是一个混沌系统，未来可能受到扰动。在未来数十亿年内，存在极小概率会有其他恒星穿越太阳系。这可能导致系统的轻微失稳，在数百万年后造成行星被逐出、相互碰撞或坠入太阳，但最有可能的结果仍是：太阳系将大体保持如今的样貌。\(^\text{[35]}\)
\begin{figure}[ht]
\centering
\includegraphics[width=8cm]{./figures/080701ee703e48af.png}
\caption{当前太阳与其在红巨星阶段达到峰值时的体积对比} \label{fig_Solar_3}
\end{figure}
太阳的主序阶段（main-sequence phase）——从开始到结束——将持续约 100 亿年，而此后所有阶段（从主序末期到白矮星前的各阶段）加起来的寿命仅约 20 亿年。\(^\text{[36]}\)太阳系将在此期间大体保持如今的状态，直到太阳核心中的氢完全转化为氦为止—— 这将发生在大约 50 亿年后。届时，太阳的主序生命将走到尽头。  

那时，太阳的核心将因引力进一步收缩，氢聚变将仅在包围惰性氦核的壳层（shell）中进行，其总能量输出将超过当前水平。太阳的外层将剧烈膨胀，直径将达到目前的约260 倍，并演化为一颗红巨星（Red Giant）。由于表面积大幅增加，其表面温度将下降，在最冷时约为2,600 K（4,220 °F），远低于主序阶段的温度。\(^\text{[36]}\)

膨胀后的太阳预计将汽化水星与金星，并使地球与火星变得不可居住（甚至可能完全摧毁地球）。\(^\text{[37][38]}\)最终，太阳核心温度将升高到足以点燃氦聚变反应；太阳将在核心燃烧氦，但这一阶段仅持续其氢燃烧阶段的极小一部分时间。由于太阳的质量不足以启动更重元素的聚变，核心中的核反应将逐渐衰竭并停止。太阳的外层将在这一过程中被抛射到太空中，留下一个致密的白矮星（White Dwarf）—— 其质量约为太阳原始质量的一半，但体积仅有地球大小。\(^\text{[36]}\)被抛射的外层气体可能形成一个行星状星云（Planetary Nebula），将太阳形成时的部分物质—— 此时已被碳等重元素所“富集”—— 重新送回星际介质（Interstellar Medium）。\(^\text{[39][40]}\)
\subsection{一般特征}
\begin{figure}[ht]
\centering
\includegraphics[width=8cm]{./figures/7b113d47a833549c.png}
\caption{这是一张经过色彩增强的从月球拍摄的太阳系多种成分的照片。左下角的三个光点自左向右依次为行星——土星、火星与水星；画面中央，太阳的日冕（corona）从月球漆黑的边缘上方升起，而月球右侧被地球反照光（earthshine）所照亮。} \label{fig_Solar_4}
\end{figure}
天文学家有时会将太阳系的结构划分为若干独立的区域。内太阳系（Inner Solar System）包括水星、金星、地球、火星以及位于小行星带中的天体；外太阳系（Outer Solar System）包括木星、土星、天王星、海王星以及位于柯伊伯带中的天体。[41]  
自从发现柯伊伯带以来，太阳系的最外层部分被视为一个独立的区域，由位于海王星之外的天体组成。[42]
\subsubsection{组成}
太阳系的主要成分是太阳，一颗G型主序星（G-type main-sequence star），其质量占整个系统已知总质量的99.86\%，在引力上完全主导着整个系统。[43]围绕太阳运行的四个最大天体——四大行星——占据剩余质量的99\%，其中仅木星和土星就合计超过90\%。太阳系的其他所有天体（包括四颗类地行星、矮行星、卫星、小行星与彗星等）总共的质量还不到太阳系总质量的0.002\%。[h]

太阳的成分大约由 **98%的氢与氦**组成，[47]  
木星与土星的成分也大体相同。[48][49]  
太阳系中存在一个由早期太阳的热量与光压形成的**成分梯度**：  
距离太阳较近、受热与光压影响更强的天体，  
主要由**熔点较高的元素**构成；  
而距离太阳较远的天体，则主要由**熔点较低的物质**组成。[50]  
在太阳系中，那条易挥发物质能够凝聚的界限称为**霜线（frost line）**，  
其位置大约在距离太阳约 **地球距离的五倍**处。[5]

---

### 轨道（Orbits）
```
